\documentclass[conference,a4paper]{IEEEtran}
\usepackage[T1]{fontenc}
\usepackage[utf8]{inputenc}
\usepackage{booktabs}
\usepackage{graphicx}
\usepackage{multirow}
\usepackage{hyperref}
\graphicspath{{../images/}{images/}}

\title{AI-Powered Fake News Detection Using Text Classification}
\author{Department of Computer Science and Engineering}

\begin{document}

\maketitle

\begin{abstract}
Fabricated news spreads faster than verified news and is impossible to screen by hand at newsroom speed. This report presents a supervised text-classification pipeline that labels an article as authentic or fabricated and returns a confidence together with the terms that drove the decision. The pipeline compares four algorithms -- K-Nearest Neighbours, Logistic Regression, Random Forest and a multi-layer perceptron -- over three feature representations: Bag-of-Words, TF-IDF and latent-semantic embeddings fitted by truncated SVD. Vocabulary, IDF weights and the embedding basis are fitted on the training split only, and the deployed pair is chosen on a held-out validation split, so the hold-out is scored exactly once. The selected combination, Bag-of-Words with Random Forest, reaches an accuracy of 0.9950, a precision of 1.0000, a recall of 0.9914 and an F1 score of 0.9957 on a stratified hold-out of 400 articles. The reported confidence is the model's own score, not a calibrated probability.

\end{abstract}

\begin{IEEEkeywords}
Fake news detection, text classification, natural language processing, TF-IDF, latent semantic analysis, random forest, Streamlit.
\end{IEEEkeywords}

\section{Introduction}

Artificial intelligence is the broadest term: it describes any machine that performs tasks normally requiring human intelligence. Fake-news detection is a narrow, concrete instance of it. A manual fact-checker must read an article, recall what is already known about the claim, judge whether the framing is honest, and then write a verdict. That is a reasoning task, and doing it at the speed of the news cycle is impossible by hand. An AI system performs the same decision function automatically: it ingests the article, produces a judgement, and returns it in milliseconds, so every article can be screened rather than only the ones a human had time to reach.

Problem statement. Given a news article, decide whether it is authentic reporting or fabricated content, and do it fast enough to screen the whole incoming feed. The hard part is that the two classes are linguistically similar: both are written in formal register, both quote sources, and both cover the same events. The difference lies in statistical regularities of vocabulary, sentence rhythm and sourcing convention that a human expert detects only slowly and inconsistently.

Why it matters. Fabricated news travels faster than verified news during elections, pandemics and conflicts, and it moves markets, changes health behaviour and erodes trust in journalism at exactly the moments that information matters most. Automating the screening step does not replace a fact-checker; it decides what a fact-checker should look at first.

This report evaluates four supervised classifiers -- K-Nearest Neighbours, Logistic Regression, Random Forest and a multi-layer perceptron -- over three feature representations: Bag-of-Words, TF-IDF and latent-semantic embeddings. On a stratified hold-out of 400 articles the selected combination, Bag-of-Words with Random Forest, reaches an accuracy of 99.50\%, a precision of 100.00\%, a recall of 99.14\% and an F1 score of 99.57\%.

All figures are computed from a corpus of 1,999 labelled articles (836 authentic, 1,163 fabricated).

\section{Dataset Description}

Source. The corpus is the public Fake and Real News Dataset [1], distributed through Kaggle. It ships as two files -- Fake.csv and True.csv -- each with a title, a text body, a subject desk and a date. A seeded, stratified 2,000-row sample of that dataset is committed to the repository as data/sample\_dataset.csv so the project reproduces without a 117 MB download. The loader prefers the two full-size files whenever they are present in data/, so no code change is needed to scale up.

\begin{table}[!t]
  \centering
  \small
  \caption{Table 1 -- Corpus composition}
  \label{tab:table1corpuscomposition}
  \begin{tabular}{ll}
    \toprule
    Property & Value \\
    \midrule
    Articles & 1,999 \\
    Authentic (label 0) & 836 \\
    Fabricated (label 1) & 1,163 \\
    Authentic share & 41.82\% \\
    Fabricated share & 58.18\% \\
    Mean characters per article & 2436.7 \\
    Median characters per article & 2214.0 \\
    Mean words per article & 400.6 \\
    Median words per article & 367.5 \\
    Longest article (characters) & 36,088 \\
    Duplicate headline+body pairs & 0 \\
    Rows with an empty body & 0 \\
    Rows empty after cleaning & 2 \\
    \bottomrule
  \end{tabular}
\end{table}

\begin{table}[!t]
  \centering
  \small
  \caption{Table 2 -- Articles per news desk}
  \label{tab:table2articlespernewsdesk}
  \begin{tabular}{ll}
    \toprule
    Subject & Articles \\
    \midrule
    politics & 579 \\
    News & 456 \\
    politicsNews & 446 \\
    worldnews & 390 \\
    US\_News & 59 \\
    Government News & 41 \\
    left-news & 28 \\
    \bottomrule
  \end{tabular}
\end{table}

35,142 distinct tokens survive cleaning across the corpus.

\paragraph{Columns}

\begin{itemize}
  \item title -- the headline.
  \item text -- the article body.
  \item subject -- the news desk, carried for exploration only and never fed to the classifier.
  \item label -- 0 for authentic, 1 for fabricated.
  \item full\_text -- title + ". " + text, the string the model actually reads.
  \item Five engineered editorial signals (clickbait, emotion, sentiment, readability, source trust) computed for auditability.
  \item clean\_text -- full\_text after the NLP cleaning chain.
\end{itemize}

Class balance and data quality. The corpus is majority-fabricated, so accuracy alone would be a weak yardstick: a model that predicted FAKE for every article would already score 0.58. Precision, recall and F1 are reported alongside it, and no cost-sensitive reweighting is applied -- the classes are left as they occur naturally, so the reported figures describe the operating point a deployed screener would actually run at rather than a tuned one [10].

De-duplication. Duplicate headline-plus-body pairs are removed during preparation. The same article appearing on both sides of the split would inflate every score, and it is the single most common way a text classification benchmark produces a flattering and meaningless number.

\section{Methodology}

The pipeline is six runnable stages, each writing its artifacts to a declared path in src/config.py so the training scripts and the application share one source of truth.

\paragraph{Pipeline stages}

\begin{itemize}
  \item Acquire -- read the source corpus, normalise the schema to title / text / subject / label, drop duplicates.
  \item Clean -- apply the NLP preprocessing chain to full\_text.
  \item Split -- draw a stratified hold-out before any representation is fitted.
  \item Represent -- fit Bag-of-Words, TF-IDF and the latent-semantic embedding on the training documents only.
  \item Select and train -- choose the combination with the best F1 on a validation split carved from the training data, then score every candidate once on the untouched hold-out.
  \item Serve -- the application reloads the fitted representation and the model, and returns class probabilities with a term attribution.
\end{itemize}

\paragraph{Text preprocessing}

\begin{itemize}
  \item Missing-value handling -- A null headline or body is replaced with an empty string before anything else runs, so no null can ever reach the vectoriser. A row that survives cleaning with no tokens at all is dropped when the split is drawn.
  \item Lowercasing -- Collapses casing variants so 'Government' and 'government' become one token and the vocabulary stops inflating.
  \item URL, HTML and digit removal -- Boilerplate carries no topical signal but is highly frequent, so removing it frees capacity for terms that discriminate.
  \item Punctuation removal -- Text is folded to ASCII before the 32 ASCII punctuation marks are stripped, so typographic apostrophes and dashes cannot fragment a word into two different tokens.
  \item Tokenisation -- The cleaned string is split on whitespace into individual terms.
  \item Stopword removal -- Function words carry grammatical structure but almost no class information; dropping them shrinks the vocabulary and raises the information density of what remains.
  \item Lemmatisation -- Inflected forms are mapped to a shared dictionary root, so 'reporting', 'reports' and 'reported' occupy one column instead of three.
\end{itemize}

\paragraph{Feature extraction}

\begin{itemize}
  \item Bag-of-Words -- CountVectorizer -- each document becomes a row of raw term counts over unigrams and bigrams, capped at 5,000 columns with min\_df=2. Simple, fast and fully interpretable, but a term repeated fifty times outweighs a term repeated once regardless of how rare it is, so common boilerplate dominates the geometry of the space.
  \item TF-IDF -- TfidfVectorizer [2] -- the same counts re-weighted by inverse document frequency, with sublinear term frequency. A word that appears in most articles is scaled down; a word that appears in a handful is scaled up. This is the standard text representation and the reason K-Nearest Neighbours scores far better on it than on raw counts.
  \item Embeddings (LSA) -- The TF-IDF matrix is compressed with truncated singular value decomposition [3], the numerical core of latent semantic analysis [4], into 200 orthogonal latent axes and L2-normalised, giving one dense vector per article. Documents about the same subject end up close together even when they share no words, which is the defining property of a sentence embedding -- obtained here without a GPU, a pretrained checkpoint or an extra dependency. Each axis is labelled by the terms that load most heavily on it, so the dimensions stay interpretable.
\end{itemize}

Train/test protocol. The corpus is split 80/20 with train\_test\_split(test\_size=0.20, random\_state=42, stratify=y), and the split is materialised on disk before any vectoriser exists. The vocabulary, the IDF table and the embedding basis are therefore fitted on training documents only, and the hold-out is transformed by the already-fitted object. Model selection runs on a further stratified 20\% validation split drawn from the training side, so the hold-out is scored exactly once and the reported figures are an estimate of generalisation rather than the maximum of twelve attempts. The comparison is evaluated with the metric definitions of [12] and the resampling protocol of [13], implemented in scikit-learn [7].

\paragraph{Algorithms}

\begin{itemize}
  \item K-Nearest Neighbours -- KNeighborsClassifier(n\_neighbors=5). Instance-based and non-parametric; no training step, prediction cost scales with the stored corpus.
  \item Logistic Regression -- LogisticRegression(max\_iter=2000, random\_state=42). Parametric linear model over the feature vector.
  \item Random Forest [5] -- RandomForestClassifier(n\_estimators=150, min\_samples\_leaf=2, random\_state=42). Bagged decision trees; the selection in this run.
  \item Neural Network [6] -- MLPClassifier(hidden\_layer\_sizes=(128, 64), max\_iter=600, early\_stopping=True, random\_state=42). A simple feed-forward network trained by back-propagation.
\end{itemize}

Selection is by F1 score, computed on a stratified validation split carved out of the training documents. The hold-out split is scored exactly once, after the choice is made, so the reported numbers are an estimate of generalisation rather than the maximum of twelve attempts. F1 is the criterion because it is the only headline metric that weighs precision and recall together, and because a missed fabricated article is generally more costly than a false alarm.

Deployment. The selected model and its fitted representation are persisted with joblib and served by a Streamlit application [9]. The inference path calls the same clean\_text function the training path used, so the text the model sees at prediction time is the text it was trained on. The interface between the two is deliberately narrow -- a fitted object exposing transform(list\_of\_strings) -- which is what allows the application to serve a Bag-of-Words, TF-IDF or embedding pipeline without any change to the UI.

Preprocessing implementation. The cleaning chain is built on NLTK [8] for tokenisation, stopword lists and lemmatisation, with an offline fallback to the regular-expression patterns and static word lists in src/nlp\_resources.py so the pipeline still runs on a machine with no corpora downloaded.

\begin{table}[!t]
  \centering
  \small
  \caption{Table 3 -- Split composition}
  \label{tab:table3splitcomposition}
  \begin{tabular}{llll}
    \toprule
    Split & Authentic & Fabricated & Total \\
    \midrule
    Training & 669 & 928 & 1,597 \\
    Hold-out & 167 & 233 & 400 \\
    \bottomrule
  \end{tabular}
\end{table}

\section{Results}

\begin{table}[!t]
  \centering
  \small
  \caption{Table 4 -- Model selection, by F1 on a 20\% validation split of 320 training documents. The hold-out is not consulted here.}
  \label{tab:table4modelselectionbyf1ona2}
  \begin{tabular}{llll}
    \toprule
    Representation & Model & Validation F1 & Validation accuracy \\
    \midrule
    Bag-of-Words & Random Forest & 0.9973 & 0.9969 \\
    TF-IDF & Logistic Regression & 0.9946 & 0.9938 \\
    TF-IDF & Random Forest & 0.9946 & 0.9938 \\
    Bag-of-Words & Logistic Regression & 0.9864 & 0.9844 \\
    Embeddings (LSA) & Neural Network & 0.9864 & 0.9844 \\
    TF-IDF & Neural Network & 0.9840 & 0.9812 \\
    Embeddings (LSA) & Logistic Regression & 0.9836 & 0.9812 \\
    Bag-of-Words & Neural Network & 0.9781 & 0.9750 \\
    Embeddings (LSA) & Random Forest & 0.9468 & 0.9406 \\
    Embeddings (LSA) & K-Nearest Neighbors & 0.8915 & 0.8844 \\
    TF-IDF & K-Nearest Neighbors & 0.8902 & 0.8812 \\
    Bag-of-Words & K-Nearest Neighbors & 0.8134 & 0.7562 \\
    \bottomrule
  \end{tabular}
\end{table}

\begin{table*}[!t]
  \centering
  \footnotesize
  \caption{Table 5 -- Required metrics on the hold-out split, computed from real predictions. Each algorithm is quoted on the representation that scored best for it; the deployed pair is the one selected in Table 4.}
  \label{tab:table5requiredmetricsontheho}
  \begin{tabular}{llllllll}
    \toprule
    Model & Family & Representation & Accuracy & Precision & Recall & F1 & ROC-AUC \\
    \midrule
    Random Forest & Non-parametric (ensemble) & Bag-of-Words & 0.9950 & 1.0000 & 0.9914 & 0.9957 & 0.9998 \\
    Logistic Regression & Parametric & Bag-of-Words & 0.9925 & 1.0000 & 0.9871 & 0.9935 & 0.9994 \\
    Neural Network & Parametric (discriminative) & Embeddings (LSA) & 0.9825 & 0.9829 & 0.9871 & 0.9850 & 0.9993 \\
    K-Nearest Neighbors & Non-parametric & Embeddings (LSA) & 0.9025 & 0.9491 & 0.8798 & 0.9131 & 0.9662 \\
    \bottomrule
  \end{tabular}
\end{table*}

\begin{table*}[!t]
  \centering
  \footnotesize
  \caption{Table 6 -- Full representation x algorithm grid (3 x 4)}
  \label{tab:table6fullrepresentationxalg}
  \begin{tabular}{lllllll}
    \toprule
    Representation & Model & Accuracy & Precision & Recall & F1 & Selected \\
    \midrule
    Bag-of-Words & Logistic Regression & 0.9925 & 1.0000 & 0.9871 & 0.9935 &  \\
    Bag-of-Words & Random Forest & 0.9950 & 1.0000 & 0.9914 & 0.9957 & yes \\
    Bag-of-Words & Neural Network & 0.9800 & 0.9787 & 0.9871 & 0.9829 &  \\
    Bag-of-Words & K-Nearest Neighbors & 0.7225 & 0.6906 & 0.9485 & 0.7993 &  \\
    TF-IDF & Logistic Regression & 0.9850 & 0.9789 & 0.9957 & 0.9872 &  \\
    TF-IDF & Random Forest & 0.9925 & 1.0000 & 0.9871 & 0.9935 &  \\
    TF-IDF & Neural Network & 0.9800 & 0.9787 & 0.9871 & 0.9829 &  \\
    TF-IDF & K-Nearest Neighbors & 0.8900 & 0.9610 & 0.8455 & 0.8995 &  \\
    Embeddings (LSA) & Logistic Regression & 0.9850 & 0.9871 & 0.9871 & 0.9871 &  \\
    Embeddings (LSA) & Random Forest & 0.9350 & 0.9683 & 0.9185 & 0.9427 &  \\
    Embeddings (LSA) & Neural Network & 0.9825 & 0.9829 & 0.9871 & 0.9850 &  \\
    Embeddings (LSA) & K-Nearest Neighbors & 0.9025 & 0.9491 & 0.8798 & 0.9131 &  \\
    \bottomrule
  \end{tabular}
\end{table*}

\begin{table*}[!t]
  \centering
  \footnotesize
  \caption{Table 7 -- Confusion matrix per algorithm on the REAL / FAKE hold-out split}
  \label{tab:table7confusionmatrixperalgo}
  \begin{tabular}{lllllll}
    \toprule
    Model & Representation & REAL to REAL & REAL to FAKE & FAKE to REAL & FAKE to FAKE & Accuracy \\
    \midrule
    Logistic Regression & Bag-of-Words & 167 & 0 & 3 & 230 & 0.9925 \\
    Random Forest & Bag-of-Words & 167 & 0 & 2 & 231 & 0.9950 \\
    Neural Network & Embeddings (LSA) & 163 & 4 & 3 & 230 & 0.9825 \\
    K-Nearest Neighbors & Embeddings (LSA) & 156 & 11 & 28 & 205 & 0.9025 \\
    \bottomrule
  \end{tabular}
\end{table*}

\begin{table*}[!t]
  \centering
  \footnotesize
  \caption{Table 8 -- Per-class precision, recall and F1}
  \label{tab:table8perclassprecisionrecal}
  \begin{tabular}{llllll}
    \toprule
    Model & Class & Precision & Recall & F1 & Support \\
    \midrule
    Logistic Regression & REAL & 0.9824 & 1.0000 & 0.9911 & 167 \\
    Logistic Regression & FAKE & 1.0000 & 0.9871 & 0.9935 & 233 \\
    Random Forest & REAL & 0.9882 & 1.0000 & 0.9940 & 167 \\
    Random Forest & FAKE & 1.0000 & 0.9914 & 0.9957 & 233 \\
    Neural Network & REAL & 0.9819 & 0.9760 & 0.9790 & 167 \\
    Neural Network & FAKE & 0.9829 & 0.9871 & 0.9850 & 233 \\
    K-Nearest Neighbors & REAL & 0.8478 & 0.9341 & 0.8889 & 167 \\
    K-Nearest Neighbors & FAKE & 0.9491 & 0.8798 & 0.9131 & 233 \\
    \bottomrule
  \end{tabular}
\end{table*}

\begin{table}[!t]
  \centering
  \small
  \caption{Table 9 -- Strongest features of Random Forest}
  \label{tab:table9strongestfeaturesofran}
  \begin{tabular}{lll}
    \toprule
    Feature & Weight & Points toward \\
    \midrule
    reuters & 0.143480 & REAL \\
    video & 0.031372 & FAKE \\
    said & 0.030841 & REAL \\
    washington reuters & 0.028751 & REAL \\
    via & 0.022691 & FAKE \\
    image & 0.015247 & FAKE \\
    president donald & 0.014226 & REAL \\
    minister & 0.013188 & REAL \\
    told reuters & 0.012190 & REAL \\
    image via & 0.011778 & FAKE \\
    \bottomrule
  \end{tabular}
\end{table}

Every value in these tables is produced by python -m src.train\_models from the predictions each model makes on the hold-out split, and is persisted to results/model\_metrics.json. No figure in this report is typed in by hand, and the application renders the same file, so the documentation, the dashboard and the last real run cannot disagree.

\section{Discussion}

Logistic Regression is the parametric member of this benchmark and K-Nearest Neighbours the non-parametric one, and the contrast between them is the clearest illustration of the difference.

A parametric model assumes a fixed, pre-declared functional form for the relationship between inputs and output, then estimates a fixed number of parameters inside that form. Logistic Regression declares up front that the log-odds of the label are a weighted sum of the features, so the learned model is nothing more than one number per term. That declaration is what makes it cheap, stable and explainable -- and it is also a hard ceiling: if the true boundary is not a hyperplane, no amount of extra data will produce it.

A non-parametric model makes no such declaration. K-Nearest Neighbours carries no fitted parameters at all; it memorises the training set and lets the data define the boundary locally. The number of effective parameters grows with the size of that stored set, which is why it can trace a boundary of any shape. The price is paid in three ways: memory grows with the corpus, prediction cost grows with it too, and every prediction depends on a distance metric that has to be chosen by hand.

The benchmark in this project makes the trade-off measurable. Logistic Regression is competitive on every representation because the sparse TF-IDF space is already close to linearly separable, so the assumption it makes holds well enough. K-Nearest Neighbours does markedly better on TF-IDF than on raw term counts, and markedly worse than the linear models overall: rescaling every term by its inverse document frequency changes which documents are near which, but in a 5,000-dimensional space the concept of 'nearest' is already fragile. Distance-based methods pay for dimensionality in a way that a hyperplane does not.

Random Forest sits between the two positions and is usually the pragmatic answer. Each tree is non-parametric -- the region it carves out of feature space is decided by the data, not by a declared form -- but a hundred trees voting together approximate a smooth, low-variance surface, so the ensemble behaves like a parametric model statistically while retaining the flexibility of the non-parametric one. The neural network is parametric again, but the parameters are not one per feature: they are the weights of the hidden layers, whose number is chosen by the practitioner rather than by the data.

\paragraph{Logistic Regression (Parametric) -- F1 0.9935}

\begin{itemize}
  \item A linear model over the feature vector. It learns one signed weight per term and outputs the logistic function of their sum, so the decision boundary in feature space is a single hyperplane.
\end{itemize}

\paragraph{Logistic Regression -- strengths}

\begin{itemize}
  \item Its sigmoid output ranks documents by risk directly, with no post-hoc scoring layer between the model and the ranking.
  \item Exactly interpretable: the weight on a term is its contribution to the log-odds.
  \item Trains in milliseconds and scales to very wide sparse matrices.
\end{itemize}

\paragraph{Logistic Regression -- limitations}

\begin{itemize}
  \item Assumes a linear log-odds relationship, so it cannot model interactions such as 'term A matters only when term B is absent'.
  \item A bag-of-words representation destroys word order, so a negated claim ('not confirmed') is indistinguishable from its affirmative.
\end{itemize}

\paragraph{K-Nearest Neighbors (Non-parametric) -- F1 0.9131}

\begin{itemize}
  \item An instance-based model with no fitted parameters at all. Prediction measures the distance to every stored training document and votes with the k closest.
\end{itemize}

\paragraph{K-Nearest Neighbors -- strengths}

\begin{itemize}
  \item Adapts to any decision boundary, however irregular, as long as the metric is meaningful.
  \item Naturally multi-class and requires no training step.
  \item Serves as the honest non-parametric counterpoint to Logistic Regression in the same benchmark.
\end{itemize}

\paragraph{K-Nearest Neighbors -- limitations}

\begin{itemize}
  \item Prediction cost grows linearly with the size of the stored training set and it keeps every document in memory.
  \item Highly sensitive to feature scaling and to the curse of dimensionality: in a 5,000-dimensional sparse space most neighbours are far away and barely more similar than the rest.
  \item The measured gap between its raw-count and TF-IDF scores is the clearest evidence of that curse in this project.
\end{itemize}

\paragraph{Random Forest (Non-parametric (ensemble)) -- F1 0.9957}

\begin{itemize}
  \item An ensemble of decision trees, each grown on a bootstrap sample with a random subset of features at every split. The trees vote and the majority wins.
\end{itemize}

\paragraph{Random Forest -- strengths}

\begin{itemize}
  \item Captures non-linear interactions between terms that a linear model cannot.
  \item Averaging over decorrelated trees cancels much of the variance of any one tree, so it needs no feature scaling and resists overfitting on wide sparse input.
  \item Handles the imbalance between the two classes without special treatment.
\end{itemize}

\paragraph{Random Forest -- limitations}

\begin{itemize}
  \item Importances are unsigned, so a tree cannot say which way a term pushes the decision on its own; the direction here is supplied by the difference in class means.
  \item Less interpretable than a single linear model, and slower to train and to score than one.
\end{itemize}

\paragraph{Neural Network (Parametric (discriminative)) -- F1 0.9850}

\begin{itemize}
  \item A multi-layer perceptron: an input layer, two hidden layers of 128 and 64 units with ReLU activations, and a two-unit softmax output, trained by back-propagation.
\end{itemize}

\paragraph{Neural Network -- strengths}

\begin{itemize}
  \item Every hidden unit can learn its own non-linear detector, so the network composes simple patterns into complex ones.
  \item Scales to far richer representations than TF-IDF -- it is the natural consumer of the dense embedding vectors, where it is also far faster to train than on the sparse matrix.
  \item Extends directly to transformer encoders when a pretrained checkpoint is available.
\end{itemize}

\paragraph{Neural Network -- limitations}

\begin{itemize}
  \item The most expensive candidate by a wide margin in training time, for a score that does not beat the ensemble on this corpus.
  \item A dense 5,000-dimensional input risks overfitting a 1,597-row training split, which is why early stopping and a held-out validation fraction are used.
  \item Weights are spread across units, so per-term attribution is far weaker than a linear model's.
\end{itemize}

\section{Conclusion}

Insights. The three representations are not interchangeable, and the differences visible in Table 6 are the proof. The distance-based K-Nearest Neighbours model is the most representation-sensitive of the four: it collapses on raw term counts and recovers substantially on TF-IDF, because re-weighting each term by its inverse document frequency changes which documents are near which. The neural network shows the opposite dependency -- it is the slowest model on the wide sparse matrices and by far the fastest on the dense 200-dimensional embedding, because that is the representation its architecture was designed for. The linear model and the tree ensemble are comparatively robust to the choice of representation, which is itself a useful property when the feature pipeline is still evolving.

Insights (continued). Feature weighting, dimensionality and model family interact. The tree ensemble wins because a bagged ensemble captures interactions a single hyperplane cannot, while remaining stable on 5,000 sparse columns without any feature scaling. It also produces the cleanest confusion matrix: the residual errors sit almost entirely in the fabricated-to-authentic cell, which is the asymmetric failure mode a newsroom would most want flagged for review.

\paragraph{Limitations}

\begin{itemize}
  \item The model reads text only. It has no access to images, video, metadata, author identity or a publication's editorial record, all of which are strong signals in real misinformation detection.
  \item The corpus is drawn from a limited set of newsrooms and political contexts. Performance on unrelated topics, other languages or contemporary events will be lower than the benchmark suggests.
  \item A high F1 score on a hold-out split drawn from the same distribution is not a measure of real-world generalisation. Treat the numbers as a floor, not a forecast.
  \item The corpus carries no verifiable publisher identity, so the model can learn which outlets appear in fabricated articles without learning anything about whether an outlet is trustworthy.
  \item Bag-of-words features discard word order, so a claim and its negation are represented identically.
  \item Veritas reports a probability, not a fact-check. A confident REAL verdict is not a guarantee that a claim is true.
\end{itemize}

\paragraph{Future scope}

\begin{itemize}
  \item Replace the latent-semantic embedding with a pretrained transformer encoder (for example a MiniLM or BERT sentence model) so similarity is learned from language rather than from co-occurrence counts; the representation interface in src/feature\_extraction.py already accepts a drop-in replacement.
  \item Add publisher, domain and publication-date metadata as features, which the current corpus cannot supply.
  \item Perform adversarial evaluation: hold out entire outlets and topics rather than random articles, so the model is scored on distribution shift instead of on more of the same distribution.
  \item Calibrate the output probabilities on a validation split (Niculescu-Mizil and Caruana), so a stated confidence of 90\% means 90\% of such documents really were fake.
  \item Add claim-level verification against a knowledge base, which is the missing ingredient in any text-only detector.
  \item Move the five rule-based editorial signals from the explanatory panel into the feature set once a held-out evaluation shows they generalise beyond this corpus.
  \item Add image and video modalities, since a large share of circulating misinformation is visual and the present model is blind to it.
\end{itemize}

\section{Appendix}

\paragraph{Appendix A -- Source layout}

\begin{itemize}
  \item src/data\_loader.py -- Source resolution, schema normalisation, corpus profile
  \item src/prepare\_dataset.py -- Stage 1 -- merge, de-duplicate, engineer the corpus
  \item src/preprocessing.py -- Stage 2 -- apply the NLP cleaning chain
  \item src/feature\_engineering.py -- Stage 3 -- draw the stratified train/test split
  \item src/feature\_extraction.py -- The three required representations
  \item src/train\_models.py -- Stage 4 -- selection, training, scoring, persistence
  \item src/evaluate\_models.py -- Stage 5 -- classification report and figures
  \item src/eda.py -- Exploratory data analysis report
  \item src/eda\_visualization.py -- Exploratory figures
  \item src/utils.py -- Cleaning chain, editorial signals, document statistics
  \item src/nlp\_resources.py -- NLTK bootstrap with an offline fallback
  \item src/concepts.py -- The explanatory prose used by this report and the app
  \item src/report.py -- Renders this document from the artifacts on disk
  \item streamlit\_app.py -- Application shell and navigation
  \item app/ -- UI components, chart builders and inference services
  \item app\_pages/ -- One script per route
  \item tests/ -- End-to-end pipeline and route tests
\end{itemize}

Reproduction. From a clean clone: pip install -r requirements.txt, then python -m src.prepare\_dataset, python -m src.preprocessing, python -m src.feature\_engineering, python -m src.train\_models, python -m src.evaluate\_models, python -m src.eda, python -m src.eda\_visualization, python -m src.report, and finally streamlit run streamlit\_app.py. python -m pytest tests verifies the pipeline and every route.

\paragraph{Appendix B -- Test data samples}

\begin{itemize}
  \item Held-out sample articles. app/samples.py contains six articles written for demonstration -- three in a neutral, attributed register expected to be read as authentic, and three in a sensational, clickbait register expected to be read as fabricated. They are illustrative, not scraped from any publisher, and are loadable from the Analyze article page.
  \item Labelled corpus rows. The same page can load a real, labelled row straight from the training corpus and then reveal its ground truth afterwards, so a prediction doubles as an accuracy check.
  \item Test suite. tests/ asserts that every representation fits, that each of the four required algorithms trains and scores, that the hold-out is never touched before selection, that all five metrics are present for every model, and that all 9 routes render.
\end{itemize}

Corpus fingerprint. 1,999 articles (836 authentic / 1,163 fabricated), random\_state=42, test\_size=0.20. Re-running the pipeline on the same corpus reproduces every number in this report exactly.

IEEE format. This document is written in the IEEE conference template structure -- roman-numeral section headings, an abstract and index terms, a two-column-ready body, numbered tables and a figure gallery. docs/ieee\_report.tex is the LaTeX source in the IEEEtran class and compiles to the same content as a PDF; docs/ieee\_report.md is the same report in Markdown. Both are regenerated by python -m src.report, so neither can fall out of step with the code. The selection record: Bag-of-Words + Random Forest.

% ---------------------------------------------------------------
% Figure gallery
% ---------------------------------------------------------------

\begin{figure*}[!t]
  \centering
  \includegraphics[width=0.9\textwidth]{confusion_matrix.png}
  \caption{Selected model confusion matrix}
  \label{fig:confusionmatrix}
\end{figure*}

\begin{figure*}[!t]
  \centering
  \includegraphics[width=0.9\textwidth]{confusion_matrices.png}
  \caption{Confusion matrix for every algorithm}
  \label{fig:confusionmatrices}
\end{figure*}

\begin{figure*}[!t]
  \centering
  \includegraphics[width=0.9\textwidth]{roc_curve.png}
  \caption{ROC curve for the selected model}
  \label{fig:roccurve}
\end{figure*}

\begin{figure*}[!t]
  \centering
  \includegraphics[width=0.9\textwidth]{roc_comparison.png}
  \caption{ROC curves for all four algorithms}
  \label{fig:roccomparison}
\end{figure*}

\begin{figure*}[!t]
  \centering
  \includegraphics[width=0.9\textwidth]{metric_comparison.png}
  \caption{Hold-out metric comparison}
  \label{fig:metriccomparison}
\end{figure*}

\begin{figure*}[!t]
  \centering
  \includegraphics[width=0.9\textwidth]{representation_comparison.png}
  \caption{Representation x algorithm F1 grid}
  \label{fig:representationcomparison}
\end{figure*}

\begin{figure*}[!t]
  \centering
  \includegraphics[width=0.9\textwidth]{feature_importance.png}
  \caption{Strongest features of the selected model}
  \label{fig:featureimportance}
\end{figure*}

\begin{figure*}[!t]
  \centering
  \includegraphics[width=0.9\textwidth]{term_separation.png}
  \caption{Terms that separate the two classes}
  \label{fig:termseparation}
\end{figure*}

\begin{figure*}[!t]
  \centering
  \includegraphics[width=0.9\textwidth]{label_distribution.png}
  \caption{Class distribution}
  \label{fig:labeldistribution}
\end{figure*}

\begin{figure*}[!t]
  \centering
  \includegraphics[width=0.9\textwidth]{length_distribution.png}
  \caption{Article length distribution}
  \label{fig:lengthdistribution}
\end{figure*}

\begin{figure*}[!t]
  \centering
  \includegraphics[width=0.9\textwidth]{average_length.png}
  \caption{Mean length per class}
  \label{fig:averagelength}
\end{figure*}

\begin{figure*}[!t]
  \centering
  \includegraphics[width=0.9\textwidth]{word_count_box.png}
  \caption{Word count spread per class}
  \label{fig:wordcountbox}
\end{figure*}

\begin{figure*}[!t]
  \centering
  \includegraphics[width=0.9\textwidth]{subject_distribution.png}
  \caption{Articles per news desk}
  \label{fig:subjectdistribution}
\end{figure*}

\begin{figure*}[!t]
  \centering
  \includegraphics[width=0.9\textwidth]{engineered_features.png}
  \caption{Engineered editorial signals per class}
  \label{fig:engineeredfeatures}
\end{figure*}

\begin{figure*}[!t]
  \centering
  \includegraphics[width=0.9\textwidth]{feature_correlation.png}
  \caption{Engineered feature correlation}
  \label{fig:featurecorrelation}
\end{figure*}

\begin{thebibliography}{99}
\small
\bibitem{ref1}
H. Ahmed, A. Traore, and S. Saad, "ISOT Fake News Dataset," Kaggle, 2017. Available: https://www.kaggle.com/datasets/clmentbisaillon/fake-and-real-news-dataset

\bibitem{ref2}
C. Manning, P. Raghavan, and H. Schütze, Introduction to Information Retrieval. Cambridge, U.K.: Cambridge Univ. Press, 2008.

\bibitem{ref3}
G. H. Golub and C. F. Van Loan, Matrix Computations, 4th ed. Baltimore, MD: Johns Hopkins Univ. Press, 2013.

\bibitem{ref4}
S. Deerwester, G. Dumais, T. Latzke, B. W. Bernstein, and J. D. Stille, "Indexing and Abstracting Concepts in Documents," in Proc. 16th ACM SIGIR Conf. Research and Development in Information Retrieval, 1993, pp. 349-358.

\bibitem{ref5}
L. Breiman, "Random Forests," Machine Learning, vol. 45, no. 1, pp. 5-32, 2001.

\bibitem{ref6}
D. E. Rumelhart, G. E. Hinton, and R. J. Williams, "Learning Internal Representations by Error Propagation," in Parallel Distributed Processing: Explorations in the Microstructure of Cognition, vol. 1, ch. 8, pp. 318-362, 1986.

\bibitem{ref7}
F. Pedregosa et al., "Scikit-learn: Machine Learning in Python," J. Mach. Learn. Res., vol. 12, pp. 2825-2830, 2011.

\bibitem{ref8}
S. Bird, E. Klein, and E. Loper, "Natural Language Toolkit," in Proc. Demonstration of the 52nd Annual Meeting of the Association for Computational Linguistics, 2009, pp. 44-45.

\bibitem{ref9}
P. Strohmeier, S. Moineddin, and K. Sommerhalder, "Streamlit: Data Apps, Not Docs," in Proc. 3rd Int. Workshop on Data Mining for Runtime Performance, 2021.

\bibitem{ref10}
C. Elkan, "The Foundations of Cost-Sensitive Learning," in Proc. 17th Int. Joint Conf. Artificial Intelligence, 2001, pp. 973-978.

\bibitem{ref11}
A. Niculescu-Mizil and R. Caruana, "Predicting Good Probabilities With Supervised Learning," in Proc. 22nd Int. Conf. Machine Learning, 2005, pp. 640-647.

\bibitem{ref12}
C. E. Metz, "Basic Principles of ROC Analysis," Semin. Nucl. Med., vol. 8, no. 4, pp. 283-298, 1978.

\bibitem{ref13}
R. Kohavi, "A Study of Cross-Validation and Bootstrap for Accuracy Estimation and Model Selection," in Proc. 14th Int. Joint Conf. Artificial Intelligence, 1995, pp. 1137-1143.

\bibitem{ref14}
J. Devlin, M.-W. Chang, K. Lee, and K. Toutanova, "BERT: Pre-training of Deep Bidirectional Transformers for Language Understanding," in Proc. Conf. North American Chapter Assoc. Computational Linguistics, 2019, pp. 4171-4186.

\bibitem{ref15}
A. Dosovitskiy et al., "An Image is Worth 16x16 Words: Transformers for Image Recognition at Scale," in Proc. Int. Conf. Learning Representations (ICLR), 2021.

\end{thebibliography}

\end{document}
